\section{Discussion}
\label{sec:discussion}

Our experiments reveal that token entropy and N-sample consistency capture orthogonal uncertainty signals for hallucination detection.

\subsection{Key Findings}

\textbf{Finding 1: Orthogonality is stronger than expected.} With only 5\% shared variance ($r = 0.228$), entropy and consistency measure fundamentally different phenomena.

\textbf{Finding 2: Consistency shows large discriminative power.} Cohen's $d = 1.068$ indicates approximately one standard deviation of separation between correct and incorrect responses.

\textbf{Finding 3: Discordant cases reveal complementary failure modes.} On 18.1\% of questions, entropy and consistency disagree---and each method achieves AUROC $> 0.75$ on its winning subset.

\subsection{Mechanistic Interpretation}

Our results support a dual-signal model of LLM uncertainty:

\begin{itemize}
    \item \textbf{Entropy captures epistemic uncertainty:} When the model lacks knowledge, its next-token distribution becomes diffuse.
    \item \textbf{Consistency captures generation stability:} When knowledge is unstable or fabricated, repeated sampling produces divergent outputs.
\end{itemize}

\subsection{Limitations}

\textbf{Single model (LLaMA-2-7B).} Our results may not generalize to other architectures or scales.

\textbf{Single benchmark (TruthfulQA).} The adversarial design may inflate effect sizes compared to naturally-occurring hallucinations.

\textbf{Hybrid detector not tested.} While we demonstrate orthogonality and complementarity, we have not built or evaluated a combined detector.

\subsection{Broader Impact}

Improved hallucination detection can reduce misinformation from LLM-generated content. Multi-signal approaches offer more robust detection than single-signal methods. We recommend combining multiple orthogonal signals to make evasion more difficult.
